\chapter{Glossary}
\label{ch:glossary}

This glossary defines recurring technical terms as they are used throughout this
thesis. Definitions are intentionally concise and reflect common usage in
information retrieval and machine learning; citations are included where a term
is tied to a particular method or formalization.

\section*{Information Retrieval}

\noindent\textbf{Information retrieval (IR).}
The study and practice of finding items relevant to an information need within a
collection. In this thesis, the items are usually documents, passages, or
document-page images, and the information need is expressed as a query.

\noindent\textbf{Corpus, document, passage, and chunk.}
A \emph{corpus} is the collection being searched. A \emph{document} is one item
in that collection, while a \emph{passage} is a span of text. In this thesis,
``passage'' and ``document'' are sometimes used interchangeably for textual
retrieval units. What counts as a document is partly a system-design choice: a
book, page, paragraph, or passage may be indexed as one item.

\emph{Chunking} divides a long document into smaller units, or \emph{chunks},
that can be indexed and retrieved independently.

\noindent\textbf{Information need, query, and relevance.}
An \emph{information need} is the underlying information a user seeks; a
\emph{query} is the expression of that need supplied to a retrieval system.
\emph{Relevance} describes how well a candidate item satisfies the information
need, not merely whether it contains the query words
\sidecite{manning2008introduction_ir}. Relevance may be binary or graded and is
task-dependent: a passage can be topically related to a query without containing
the evidence needed to answer it.

\noindent\textbf{Indexing, retrieval, and reranking.}
\emph{Indexing} preprocesses a corpus into a data structure that supports fast
search. \emph{Retrieval} uses that index to obtain an initial set of candidates.
\emph{Reranking} then applies a typically slower and more expressive model to
reorder a limited candidate set.

\noindent\textbf{Sparse, lexical, and dense retrieval.}
\emph{Lexical retrieval} primarily matches explicit terms shared by queries and
documents. Its representations are usually \emph{sparse}: most coordinates,
often associated with vocabulary terms, are zero. \emph{Dense retrieval} uses
learned, low-dimensional vectors whose coordinates are generally non-zero and
can capture semantic similarity even when the query and document use different
words.

\noindent\textbf{Inverted index.}
A data structure composed of a dictionary of terms and, for each term, a
\emph{postings list} identifying the documents, and often positions, in which it
occurs. It is the standard basis for efficient lexical retrieval at corpus scale
\sidecite{manning2008introduction_ir}.

\noindent\textbf{Nearest-neighbor search and vector index.}
Given a query vector, nearest-neighbor search finds the indexed vectors that are
closest under a chosen distance or similarity function. Exact search compares
the query with every vector; approximate nearest-neighbor (ANN) methods use
specialized data structures to trade a small amount of recall for much lower
latency and memory use at scale. A \emph{vector index} stores embeddings and
supports this search efficiently \sidecite{johnson2017faiss}.

\noindent\textbf{Term frequency--inverse document frequency (TF--IDF).}
A family of term-weighting schemes that gives a term more weight when it occurs
frequently in a document but less weight when it occurs in many documents
\sidecite{sparck_jones_statistical_1972}. For a corpus
$\mathcal{D}=\{d_1,\ldots,d_N\}$, a common form is
\[
  \operatorname{tfidf}(t,d,\mathcal{D})
  =
  \operatorname{tf}(t,d)\operatorname{idf}(t,\mathcal{D}),
  \qquad
  \operatorname{idf}(t,\mathcal{D})
  =
  \log\!\left(\frac{N}{\operatorname{df}(t)}\right),
\]
where $\operatorname{tf}(t,d)$ measures the frequency of term $t$ in document
$d$, and $\operatorname{df}(t)$ is the number of corpus documents containing
$t$. Many implementations use logarithmic scaling or smoothing. In the classic
vector-space model, TF--IDF weights form query and document vectors that can be
ranked using a similarity function such as cosine similarity
\sidecite{manning2008introduction_ir}.

\noindent\textbf{Okapi BM25.}
A widely used lexical ranking function that combines inverse document
frequency, saturating term frequency, and document-length normalization
\sidecite{robertson1994bm25,robertson_okapi_1994}. Using the notation introduced
in the history chapter, its score is
\[
  \operatorname{BM25}(q,d)
  =
  \sum_{t\in q}
  \operatorname{idf}(t)
  \frac{f_{t,d}(k_1+1)}
       {f_{t,d}+k_1\!\left(1-b+b\,\frac{|d|}{\operatorname{avgdl}}\right)},
\]
with
\[
  \operatorname{idf}(t)
  =
  \log\!\left(
  \frac{N-\operatorname{df}(t)+0.5}
       {\operatorname{df}(t)+0.5}
  \right).
\]
Here $f_{t,d}$ is the count of $t$ in $d$, $|d|$ is the document length,
$\operatorname{avgdl}$ is the corpus-average document length, $k_1>0$ controls
term-frequency saturation, and $b\in[0,1]$ controls length normalization.

\section*{Representations and Matching}

\noindent\textbf{Embedding.}
A learned numerical representation of an item as a vector or set of vectors.
Text embeddings may represent tokens, passages, queries, or entire documents;
visual embeddings may represent image patches or complete images. Retrieval
models are trained so that geometrical relationships between embeddings are
useful for matching.

\noindent\textbf{Embedding space and latent space.}
An \emph{embedding space} is the vector space in which embeddings are compared.
A \emph{latent space} is a more general learned representation space whose
coordinates are not directly observed or manually assigned. In this thesis,
``latent space'' and ``representation space'' often refer to the internal
learned geometry of representations exposed for downstream comparison or retrieval.

\noindent\textbf{Similarity function.}
A function that assigns a matching score to two representations. Common choices
include the dot product and cosine similarity. Cosine similarity compares the
angle between vectors,
\[
  \cos(\mathbf{x},\mathbf{y})
  =
  \frac{\mathbf{x}^{\mathsf{T}}\mathbf{y}}
       {\lVert\mathbf{x}\rVert_2\lVert\mathbf{y}\rVert_2},
\]
and therefore ignores their overall magnitudes.

\noindent\textbf{Single-vector and multi-vector representation.}
A \emph{single-vector} model compresses each query or document into one vector.
A \emph{multi-vector} model retains several vectors, commonly one per token or
image patch, allowing a more fine-grained matching function at the cost of a
larger index.

\noindent\textbf{Pooling.}
An operation that aggregates several representations into fewer
representations, often a single vector. \emph{Mean pooling} averages vectors,
\emph{max pooling} retains the maximum according to a chosen aggregation
criterion, and
\emph{special-token pooling} uses the representation of a designated token such
as \texttt{[CLS]} or the final token. Token pooling can also group nearby or
redundant token vectors to reduce the storage and computation required by
multi-vector retrieval. In MaxSim late interaction, for example, the maximum is
taken over document-token similarity scores for each query token, rather than
over embedding dimensions.

\noindent\textbf{Contextual embedding.}
An embedding whose value depends on surrounding input rather than representing
an item in isolation. A contextual token embedding changes when the token's
sentence changes; similarly, a contextual chunk embedding may incorporate
information from other passages in the same document or from the wider corpus.

\noindent\textbf{Bi-encoder (dual encoder).}
A model that encodes the query and document independently, often into
single-vector embeddings. Document representations can therefore be computed
offline and indexed, while retrieval reduces to efficient nearest-neighbor
search. This architecture favors scalability but limits direct query--document
interaction \sidecite{reimers2019sbert}.

\noindent\textbf{Cross-encoder.}
A model that processes a query and candidate document jointly, allowing
self-attention across all query and document tokens before producing a relevance
score. Cross-encoders are expressive but cannot precompute a query-independent
document score, so they are generally used for reranking rather than exhaustive
corpus search.

\noindent\textbf{Late interaction.}
A retrieval architecture between bi-encoders and cross-encoders. Queries and
documents are encoded independently into multi-vector representations, but
their token- or patch-level vectors interact at scoring time. ColBERT's common
\emph{MaxSim} operator finds, for every query vector, the most similar document
vector and sums these maxima \sidecite{khattab2020colbert}. This preserves
offline document indexing while enabling finer matching than a single-vector
similarity.

\noindent\textbf{Contrastive learning.}
A training approach that makes representations of matched or relevant examples
more similar and representations of unmatched examples less similar. In
retrieval, a query and its relevant document form a positive pair, while other
documents serve as negatives; \emph{hard negatives} are non-relevant candidates
that are difficult for the current model to distinguish from positives.

\noindent\textbf{Positive, negative, and hard-negative examples.}
A \emph{positive} is an example that should match the anchor input, such as a
relevant document for a query. A \emph{negative} should not match it.
\emph{Hard negatives} are non-relevant examples that nevertheless appear
plausible to the model, for example because they are topically similar or
highly ranked by another retriever. They generally provide a stronger training
signal than randomly sampled negatives, but false negatives can teach the model
to separate genuinely relevant pairs.

\noindent\textbf{Contrastive loss.}
An objective used to implement contrastive learning. It increases the score of
positive pairs relative to negative pairs, commonly by applying a softmax and
cross-entropy over the candidates in a batch. A temperature parameter may
control how sharply score differences affect the loss.

\noindent\textbf{Late chunking.}
A method that encodes a long document before pooling token representations into
individual chunk embeddings. Unlike independently encoding each chunk, this
allows a chunk representation to incorporate information from surrounding
document context \sidecite{LateChunking}.

\section*{Model Architectures}

\noindent\textbf{Token and tokenizer.}
A \emph{token} is a discrete unit processed by a language model, such as a word,
subword, character, punctuation mark, or special symbol. A \emph{tokenizer}
defines the vocabulary and the procedure that converts raw text into tokens and
back into text.

\noindent\textbf{Pretraining, fine-tuning, and inference.}
\emph{Pretraining} learns general-purpose model parameters from a large and
broad dataset, usually before a specific downstream task is fixed.
\emph{Fine-tuning} continues training those parameters, or a smaller set of
added parameters, on task- or domain-specific data. \emph{Inference} is the use
of the trained model to encode, score, classify, or generate outputs without
updating its parameters.

\noindent\textbf{Transformer.}
A neural-network architecture built around attention mechanisms rather than
recurrence. Transformer layers use self-attention to let each token combine
information from other tokens, followed by position-wise feed-forward
transformations \sidecite{vaswani_attention_2017}. Transformers can be organized
as encoders, decoders, or encoder--decoder models and underpin most language and
vision models discussed in this thesis.

\noindent\textbf{Attention, self-attention, and causal attention.}
\emph{Attention} computes a weighted combination of value vectors according to
the compatibility of queries and keys. In \emph{self-attention}, queries, keys,
and values are derived from the same sequence. \emph{Bidirectional} self-attention
allows a token to use tokens on both sides, whereas \emph{causal} attention masks
future positions so that a token can depend only on earlier tokens.

\noindent\textbf{Encoder and decoder.}
An \emph{encoder} maps an input into contextual representations and commonly
uses bidirectional attention. An autoregressive \emph{decoder} predicts a
sequence one token at a time using causal attention. An
\emph{encoder--decoder} model first encodes an input and then generates an
output conditioned on those encoded representations.

\noindent\textbf{Autoregressive model.}
A model that factorizes the probability of a sequence into successive
next-token predictions. During generation, each new token is sampled or selected
conditioned on the preceding tokens and is then appended to the context used to
predict the next one.

\noindent\textbf{Context window.}
The maximum amount of input, measured in tokens, that a model can process in one
forward pass. Material outside this window must be discarded, compressed, or
handled through mechanisms such as chunking or retrieval. A long context window
increases capacity to process complete documents but also raises computation
and memory costs.

\noindent\textbf{Large and small language models (LLMs and SLMs).}
A language model assigns probabilities to token sequences and can generate text.
Following the operational definition used in the CroissantLLM work, an LLM is
pretrained on a large text corpus, supports text generation, and enables
transfer through fine-tuning or prompting \sidecite{rogers2024positionkeyclaimsllm}.
The cited definition proposes one billion pretraining tokens as a lower
threshold; under these criteria, CroissantLLM is an LLM despite its relatively
small parameter count compared with current frontier models.
``Small language model'' is a relative term for a lower-parameter language model
designed, for example, for lower latency, memory use, or local inference; there
is no universal parameter threshold separating SLMs from LLMs.

\noindent\textbf{Vision Transformer (ViT) and image patch.}
A Vision Transformer divides an image into fixed-size regions called
\emph{patches}, projects them into vectors, and processes the resulting sequence
with a transformer \sidecite{dosovitskiy_image_2020}. A patch embedding is
therefore analogous to a text-token embedding, although it represents a spatial
image region.

\noindent\textbf{Vision--language model (VLM).}
A model that processes both visual and textual inputs in a shared architecture
or aligned representation spaces. VLMs may be generative, contrastive, or
encoder-based; in this thesis they are frequently adapted to produce embeddings
for image--text retrieval.

\noindent\textbf{Modality alignment.}
Training that makes representations from different modalities, such as images
and text, mutually compatible. Alignment can map paired inputs into nearby
regions of a shared embedding space or teach a multimodal model to process one
modality through representations expected by another. Contrastive image--text
training is a common alignment method \sidecite{radford_learning_2021}.

\noindent\textbf{Early fusion and dual-encoder multimodal models.}
An \emph{early-fusion} model combines representations from different modalities
inside a shared backbone, allowing cross-modal interaction throughout much of
the network. A multimodal \emph{dual encoder} processes each modality with a
separate encoder and compares only their final representations. Dual encoders
support efficient precomputation, while early fusion enables richer interaction
at greater computational cost.

\noindent\textbf{Quantization.}
The representation of model parameters, activations, or stored embeddings with
lower-precision values than those used in standard floating-point computation.
Quantization reduces memory, storage, and often latency, but aggressive
compression can reduce model or retrieval quality.

\section*{Document AI and Evaluation}

\noindent\textbf{Optical character recognition (OCR).}
The process of converting text visible in an image or scanned document into
machine-readable characters. OCR-based retrieval first extracts text, whereas
visual document retrieval can index page images directly.

\noindent\textbf{Visual document retrieval (VDR).}
Retrieval in which documents, usually pages, are represented from their visual
appearance rather than exclusively from parsed text. This retains evidence
carried by layout, typography, tables, figures, and spatial relationships, as
well as by visible words.

\noindent\textbf{Retrieval-augmented generation (RAG).}
A system pattern in which a retriever selects external evidence and a generative
model produces an output conditioned on that evidence. The original RAG
formulation combined a sequence-to-sequence generator with a learned dense
Wikipedia retriever \sidecite{lewis_retrieval-augmented_2020}; current usage also
covers pipelines with lexical or hybrid retrieval, reranking, query rewriting,
and iterative search.

\noindent\textbf{Parametric and non-parametric knowledge.}
\emph{Parametric knowledge} is information encoded implicitly in a model's
learned weights. \emph{Non-parametric knowledge} is stored in an external
resource, such as a document corpus or vector index, that can be searched or
updated without retraining the model. RAG combines a parametric generator with
access to non-parametric knowledge \sidecite{lewis_retrieval-augmented_2020}.

\noindent\textbf{Test collection and relevance judgment.}
A standard retrieval test collection contains a document corpus, a set of
information needs represented as queries, and judgments identifying which
documents are relevant to each need \sidecite{manning2008introduction_ir}.
Retrieval metrics summarize performance with respect to these judgments, which
act as the evaluation ground truth.

\noindent\textbf{Recall@$k$.}
The fraction of all relevant items for a query that occur among the first $k$
retrieved results. When a benchmark contains one known relevant item per query,
Recall@$k$ is equivalently the fraction of queries for which that item appears
in the top $k$.

\noindent\textbf{Precision@$k$.}
The fraction of the first $k$ retrieved items that are relevant. Unlike
Recall@$k$, it does not account for relevant items that exist elsewhere in the
corpus.

\noindent\textbf{Mean reciprocal rank (MRR).}
For each query, reciprocal rank is the inverse of the rank of the first relevant
result, or zero if no relevant result is retrieved. MRR averages this value over
queries and therefore emphasizes placing at least one relevant result near the
top of the ranking.

\noindent\textbf{Average precision and mean average precision (AP and MAP).}
Average precision summarizes a ranking by averaging precision at the ranks
where relevant items occur. Mean average precision averages AP over queries. It
rewards retrieving all relevant items and placing them early, and is most
directly applicable to binary relevance judgments.

\noindent\textbf{Normalized discounted cumulative gain (nDCG@$k$).}
A ranking metric for graded relevance. It accumulates relevance gains through
rank $k$, discounts gains appearing lower in the ranking, and divides by the
score of an ideal ordering. The normalized score lies between zero and one,
with one indicating an ideal ranking under the supplied relevance judgments.

\section*{Language-Model Training and Evaluation}

\noindent\textbf{Perplexity.}
An exponentiated average negative log-likelihood that measures how much
probability a language model assigns to a token sequence. Lower perplexity means
the model predicts the evaluated token sequence with greater confidence.
Perplexities are not directly comparable across models that use different
tokenizers or evaluation protocols.

\noindent\textbf{Scaling law.}
An empirical relationship describing how model loss or downstream performance
changes as a power-law-like function of factors such as parameter count,
training data, or compute. Scaling laws can guide how a fixed training budget is
allocated, but their predictions depend on the architecture, data distribution,
and range over which they were fitted
\sidecite{scaling-laws-lm,hoffmann2022training}.

\noindent\textbf{Zero-shot and few-shot learning.}
In a \emph{zero-shot} evaluation, a model performs a task without task-specific
demonstrations in its input. In a \emph{few-shot} evaluation, the prompt
contains a small number of example input--output pairs. These terms describe
the evaluation context and do not imply that the model has never encountered
related tasks or data during training.

\section*{CroissantLLM-Specific Terms}

\noindent\textbf{Effective capacity ratio.}
The \emph{effective parameters} of a multilingual model for one language are the
number of non-embedding parameters a comparable monolingual model would require
to match its performance in that language \sidecite{pmlr-v202-fernandes23a}.
The effective capacity ratio divides this estimate by the multilingual model's
number of non-embedding parameters. It measures how much of the multilingual
model's capacity appears effectively available to the language in question. In
the joint scaling-law experiments reported for CroissantLLM, models trained on
equal proportions of English and French reach an effective capacity ratio above
82\% for French, suggesting substantial capacity sharing between the two
languages.

\noindent\textbf{Tokenizer fertility.}
The average number of subword tokens produced per tokenized word
\sidecite{rust2021good}. Lower fertility means that a tokenizer represents the
same text with fewer tokens; its theoretical lower bound is one token per word
under a word-based measurement.
